% from ega0-1-fr.tex
\oldpage[0\textsubscript{I}]{20}

si on identifie $M_f$ et $M_{fg}$ à $\varinjlim M_{f,n}$ et
$\varinjlim M_{fg,n}$ respectivement, $\rho_{fg,f}$ s'identifie à la
\emph{limite inductive} des applications
$\rho_{fg,f}^n:M_{f,n}\to M_{fg,n}$ définies par
$\rho_{fg,f}^n(z)=g^nz$. En effet, cela résulte immédiatement de la
commutativité du diagramme
\[
\begin{tikzcd}[column sep=huge,row sep=large]
M_{f,n} \arrow[r,"\rho_{fg,f}^n"] \arrow[d,"\varphi_n^f"']
  & M_{fg,n} \arrow[d,"\varphi_n^{fg}"] \\
M_f \arrow[r,"\rho_{fg,f}"]
  & M_{fg}.
\end{tikzcd}
\]
\end{env}

\subsection{Support d'un module}
\label{subsection:0.1.7-fr}

\begin{env}[1.7.1]
\label{0.1.7.1-fr}
Étant donné un $A$-module $M$, on appelle \emph{support} de $M$ et on note
$\operatorname{Supp}(M)$ l'ensemble des idéaux premiers $\mathfrak p$ de $A$
tels que $M_{\mathfrak p}\neq0$. Pour que $M=0$, il faut et il suffit que
$\operatorname{Supp}(M)=\varnothing$, car si $M_{\mathfrak p}=0$ pour tout
$\mathfrak p$, l'annulateur d'un élément $x\in M$ ne peut être contenu dans
aucun idéal premier de $A$, donc est $A$ tout entier.
\end{env}

\begin{env}[1.7.2]
\label{0.1.7.2-fr}
Si $0\to N\to M\to P\to0$ est une suite exacte de $A$-modules, on a
\[
  \operatorname{Supp}(M)=\operatorname{Supp}(N)\cup\operatorname{Supp}(P)
\]
car pour tout idéal premier $\mathfrak p$ de $A$, la suite
$0\to N_{\mathfrak p}\to M_{\mathfrak p}\to P_{\mathfrak p}\to0$ est exacte
(1.3.2) et pour que $M_{\mathfrak p}=0$, il faut et il suffit que
$N_{\mathfrak p}=P_{\mathfrak p}=0$.
\end{env}

\begin{env}[1.7.3]
\label{0.1.7.3-fr}
Si $M$ est somme d'une famille $(M_\lambda)$ de sous-modules,
$M_{\mathfrak p}$ est somme des $(M_\lambda)_{\mathfrak p}$ pour tout idéal
premier $\mathfrak p$ de $A$ (1.3.3 et 1.3.2), donc
\[
  \operatorname{Supp}(M)=\bigcup_\lambda\operatorname{Supp}(M_\lambda).
\]
\end{env}

\begin{env}[1.7.4]
\label{0.1.7.4-fr}
Si $M$ est un $A$-module \emph{de type fini}, $\operatorname{Supp}(M)$ est
l'ensemble des idéaux premiers \emph{contenant l'annulateur de $M$}. En effet,
si $M$ est monogène et engendré par $x$, dire que $M_{\mathfrak p}=0$ signifie
qu'il existe $s\notin\mathfrak p$ tel que $s.x=0$, donc que $\mathfrak p$ ne
contient pas l'annulateur de $x$. Si maintenant $M$ admet un système fini
$(x_i)_{1\leq i\leq n}$ de générateurs et si $\mathfrak a_i$ est l'annulateur
de $x_i$, il résulte de (1.7.3) que $\operatorname{Supp}(M)$ est l'ensemble des
$\mathfrak p$ contenant l'un des $\mathfrak a_i$, ou, ce qui revient au même,
l'ensemble des $\mathfrak p$ contenant $\mathfrak a=\bigcap_i\mathfrak a_i$,
qui est l'annulateur de $M$.
\end{env}

\begin{env}[1.7.5]
\label{0.1.7.5-fr}
Si $M$ et $N$ sont deux $A$-modules \emph{de type fini}, on a
\[
  \operatorname{Supp}(M\otimes_A N)
  =\operatorname{Supp}(M)\cap\operatorname{Supp}(N).
\]
Il s'agit de voir que si $\mathfrak p$ est un idéal premier de $A$, la
condition $M_{\mathfrak p}\otimes_{A_{\mathfrak p}}N_{\mathfrak p}\neq0$ est
équivalente à «~$M_{\mathfrak p}\neq0$ et $N_{\mathfrak p}\neq0$~» (compte
tenu de (1.3.4)). Autrement dit, il s'agit de voir que si $P$, $Q$ sont deux
modules de type fini sur un anneau \emph{local} $B$, non réduits à $0$, alors
$P\otimes_B Q\neq0$. Soit $\mathfrak m$ l'idéal maximal de $B$. En vertu du
lemme de Nakayama, les espaces vectoriels $P/\mathfrak mP$ et
$Q/\mathfrak mQ$ ne sont pas réduits à $0$, donc il en est de même de leur
produit tensoriel
\[
  (P/\mathfrak mP)\otimes_{B/\mathfrak m}(Q/\mathfrak mQ)
  =(P\otimes_B Q)\otimes_B(B/\mathfrak m),
\]
d'où la conclusion.

En particulier, si $M$ est un $A$-module de type fini, $\mathfrak a$ un idéal
de $A$, $\operatorname{Supp}(M/\mathfrak aM)$ est l'ensemble des idéaux
premiers contenant à la fois $\mathfrak a$ et l'annulateur $\mathfrak n$ de
$M$ (1.7.4), autrement dit l'ensemble des idéaux premiers contenant
$\mathfrak a+\mathfrak n$.
\end{env}

